\documentclass[11pt]{article}
\usepackage[margin=1in]{geometry}
\usepackage{amsmath,amssymb,amsthm}
\usepackage[T1]{fontenc}
\usepackage{lmodern}
\usepackage[colorlinks=true,linkcolor=blue,citecolor=blue,urlcolor=blue]{hyperref}
\newtheorem{theorem}{Theorem}[section]
\newtheorem{lemma}[theorem]{Lemma}
\newtheorem{proposition}[theorem]{Proposition}
\newtheorem{corollary}[theorem]{Corollary}
\theoremstyle{remark}
\newtheorem{remark}[theorem]{Remark}
\newcommand{\Prob}{\mathbb P}
\newcommand{\E}{\mathbb E}
\newcommand{\D}{\mathcal D}
\newcommand{\M}{\mathcal M}
\newcommand{\Q}{\mathbb Q}
\title{Tail bounds for Markov equivalence classes\protect\\of uniform random DAGs}
\author{Kaiwen Jin\thanks{Contact: \href{mailto:yui.ui.siki@gmail.com}{yui.ui.siki@gmail.com}.}}
\date{October 13, 2026}
\begin{document}
\maketitle
\begin{abstract}
Let $G_n$ be uniform among the simple directed acyclic graphs on $n$ labeled
vertices. We prove that its Markov equivalence class satisfies
\[
\Prob\bigl(|\M(G_n)|>t\bigr)
 \le C\exp\left[-c\left(\frac{\log t}{\log\log t}\right)^2\right]
\]
for all sufficiently large $t$, uniformly in $n$. This implies the upper tail
with exponent $2-\varepsilon$ asked for by Schmid and Sly. The main estimate is
a Gaussian tail for the number of vertices incident to reversible edges.
Within a large essential-graph component, absence of unshielded colliders is
costly; between many small components, their common-parent condition supplies
the remaining quadratic cost. Source-layer enumeration transfers these
estimates from independent forward edges to the uniform-DAG measure.
Complete skeletons show that the logarithmic scale of the bound is sharp up to
constants in the exponent.
\end{abstract}

\section{Introduction and results}
Let $\D_n$ be the set of simple directed acyclic graphs (DAGs) on
$[n]=\{1,\ldots,n\}$, and let $G_n$ be uniform on $\D_n$. The skeleton of a
DAG forgets its edge orientations. An unshielded collider, or $v$-structure,
is a triple $u\to v\leftarrow w$ with $u,w$ nonadjacent. We write
\[
\M(G)=\{H\in\D_n:H\text{ has the same skeleton and unshielded colliders as }G\}.
\]
These are the Markov equivalent DAGs; see \cite[Theorem 2.1]{AMP}.
All logarithms below are natural.

Schmid and Sly \cite[Theorem 1.1]{SS} prove a uniform upper tail with power
$1+1/20$ on $\log t$ and ask, immediately after that theorem, whether the
power can approach two from below. We prove the following stronger bound.
The distribution here is uniform over labeled DAGs. Choosing a uniform
vertex order and then independent forward edges gives a different law.

\begin{theorem}\label{thm:main}
There are universal constants $c,C>0$ and $t_0>e$ such that, for every integer
$n\ge1$ and every real $t\ge t_0$,
\begin{equation}\label{eq:main}
 \Prob\bigl(|\M(G_n)|>t\bigr)
 \le C\exp\left[-c\left(\frac{\log t}{\log\log t}\right)^2\right].
\end{equation}
\end{theorem}

\begin{corollary}\label{cor:target}
For every $\varepsilon\in(0,1)$ there are $C_\varepsilon<\infty$ and
$t_\varepsilon>1$ such that, simultaneously for every integer $n\ge1$ and every
real $t\ge t_\varepsilon$,
\[
 \Prob\bigl(|\M(G_n)|>t\bigr)
 \le C_\varepsilon\exp[-(\log t)^{2-\varepsilon}].
\]
\end{corollary}

The essential graph of $G$ retains an arrow when all members of $\M(G)$
give it the same orientation, and replaces every other edge by an undirected
edge. Let $C_1,\ldots,C_k$ be its undirected connected components with at
least two vertices, and set
\[
 R(G)=\sum_{i=1}^k |C_i|,
\]
with $R(G)=0$ when $k=0$. Thus $R$ counts the vertices incident to reversible
edges. Our main combinatorial estimate is
\begin{theorem}\label{thm:vertices}
There are universal constants $c,C>0$ such that, for all integers $n,m\ge1$,
\begin{equation}\label{eq:vertices}
 \Prob(R(G_n)\ge m)\le C e^{-cm^2}.
\end{equation}
\end{theorem}
The factorial bound $|\M(G)|\le R(G)!$ then gives Theorem~\ref{thm:main}.
The proof uses the source-layer enumeration underlying the tower method of
Schmid and Sly \cite[Section 2, Lemma 3.1 and Proposition 3.2]{SS}, with
elementary bounds sufficient for our purpose. We give those bounds in full.
The additional quadratic estimate uses both the internal structure of each
essential component and the arrows between components.

\section{Necessary conditions on essential components}
We use the essential-graph characterization of Andersson, Madigan and
Perlman \cite[Propositions 4.1--4.2 and Theorem 4.1]{AMP}. The essential graph
is a chain graph: it has no semidirected cycle, meaning a cycle containing at
least one arrow whose arrows all point consistently around the cycle. Its
undirected components are chordal, and each member of the equivalence class
orients each component perfectly, that is, acyclically and without an
unshielded collider inside the component. The essential graph also has no
induced configuration $u\to v\text{---}w$ with $u,w$ nonadjacent.
The quoted locations are printed pages 10--11 of the primary-author technical
report \cite{AMP}; perfect orientations are defined in its Appendix A.

\begin{lemma}\label{lem:components}
Every induced graph $G[C_i]$ has no unshielded collider. The directed graph
obtained by contracting all essential components is acyclic. If $u\notin C_i$
and $u\to v$ for some $v\in C_i$, then $u\to w$ for every $w\in C_i$.
Moreover,
\begin{equation}\label{eq:factorial}
 |\M(G)|\le\prod_{i=1}^k |C_i|!\le R(G)!.
\end{equation}
Every reversible edge is absent from all unshielded colliders of $G$.
\end{lemma}
\begin{proof}
The internal assertion follows from perfect orientation. A directed cycle
after contraction would yield a semidirected cycle before contraction, using
undirected paths inside the components.
For the parent assertion, take an undirected edge $v\text{---}w$ within $C_i$.
The absence of an induced $u\to v\text{---}w$ forces $u,w$ to be adjacent.
Their edge is directed in the essential graph because $u$ is outside the
component. It cannot point from $w$ to $u$, since
$u\to v\text{---}w\to u$ would be a semidirected cycle. Hence $u\to w$.
Propagation along undirected paths proves the assertion for every $w$.

All arrows between components are fixed throughout the class. An acyclic
orientation of a component is determined by some topological order, and each
order determines at most one orientation of its fixed skeleton. There are at
most $|C_i|!$ such orientations. The product bounds the class size; the second
inequality in \eqref{eq:factorial} is the multinomial inequality. When $R=0$
both sides equal one. Finally, either arrow of a specified unshielded collider
must keep its direction throughout the class, so it cannot be reversible.
\end{proof}

\section{Source layers and comparison with independent edges}
Delete all sources of a DAG, then all sources of the remaining graph, and
continue. Write the resulting nonempty layers as $H_1,\ldots,H_T$, with
$h_i=|H_i|$, and put $h_i=0$ for $i>T$. The process terminates because every
nonempty finite DAG has a source. Edges point from an earlier to a later
layer, and every vertex of $H_{i+1}$ has at least one parent in $H_i$.
Conversely these two conditions characterize the source layers.

Conditionally on the graph, order each layer uniformly and independently,
and concatenate the orders. We call the positions in this list ranks.
This auxiliary randomness does not change the law of $G_n$. Let $F_q$ be
the induced graph on its first $q$ ranks. For comparison, $\Q_q$ denotes
the law on $[q]$ which includes each possible arrow $a\to b$, $a<b$,
independently with probability $1/2$. An exact forward graph has $\Q_q$-mass
$2^{-\binom q2}$.

Put $a_n=|\D_n|$, $a_0=1$, and
\[
 B_n=\frac{a_n}{n!2^{\binom n2}}.
\]
For a positive composition $h=(h_1,\ldots,h_T)$ of $n$, define
\begin{equation}\label{eq:weight}
 W(h)=\prod_{i=1}^T\frac{2^{-\binom{h_i}{2}}}{h_i!}
       \prod_{i<T}(1-2^{-h_i})^{h_{i+1}}.
\end{equation}
The empty composition has weight one.

\begin{lemma}\label{lem:enumeration}
We have $B_n=\sum_{|h|=n}W(h)$ and
$\Prob(h(G_n)=h)=W(h)/B_n$. Conditionally on the ordered labeled layers,
all edges joining layers at distance at least two are mutually independent
fair bits. This remains true after the intralayer orders are chosen.
\end{lemma}
\begin{proof}
The labeled layer sets can be chosen in $n!/\prod_i h_i!$ ways.
Adjacent layers have $(2^{h_i}-1)^{h_{i+1}}$ possible edge sets, since
each vertex of the later layer needs a parent in the earlier one.
Nonadjacent layers have unrestricted edges. The total is therefore
\[
 \frac{n!}{\prod_i h_i!}\,2^{\sum_{i<j}h_i h_j}
 \prod_{i<T}(1-2^{-h_i})^{h_{i+1}}.
\]
Divide by $n!2^{\binom n2}$ and use
$\binom n2=\sum_i\binom{h_i}{2}+\sum_{i<j}h_i h_j$.
No adjacent-layer restriction involves a nonadjacent-layer edge, which also
proves independence. The auxiliary orders are chosen independently of these
edge bits.
\end{proof}

\begin{lemma}\label{lem:layer}
For every $n\ge1$, every layer index $i\ge1$, and every integer $x\ge32$,
\begin{equation}\label{eq:layer}
 \Prob(h_i\ge x)\le4\,2^{-x^2/4}.
\end{equation}
\end{lemma}
\begin{proof}
Fix $y\ge x$. Replace an entry $h_i=y$ by $y$ consecutive ones, obtaining
a composition $\widetilde h$ of the same $n$. For fixed $i,y$ this map is
injective, since the specified singleton entries can be collapsed again.
Let $s=h_{i+1}$, taking $s=0$ at the last layer. The preceding interface,
if present, contributes a ratio at most one. The new singleton interfaces
contribute $2^{y-1}$ to $W(h)/W(\widetilde h)$, and the following interface
contributes at most $2^s$. Thus
\[
 \frac{W(h)}{W(\widetilde h)}
 \le\frac{2^{-\binom{y-1}{2}+s}}{y!}.
\]
If $s\le x^2/16$, then $y\ge x\ge32$ gives
\[
 -\binom{y-1}{2}+s
 \le-\frac{y^2}{2}+\frac{3y}{2}-1+\frac{y^2}{16}
 \le-\frac{y^2}{4}.
\]
Summing over the injective image and dividing by $B_n$ yields
$\Prob(h_i=y,h_{i+1}\le x^2/16)\le2^{-y^2/4}$.
Consecutive terms in the sum over $y\ge x$ have ratio less than $1/2$, so
\[
 \Prob(h_i\ge x)
 \le\Prob\left(h_{i+1}\ge\lfloor x^2/16\rfloor+1\right)
       +2\,2^{-x^2/4}.
\]
Iterate with $x_0=x$ and $x_{j+1}=\lfloor x_j^2/16\rfloor+1$.
Then $x_{j+1}\ge2x_j$ for $x_j\ge32$. Once $x_j>n$ the remaining
probability vanishes. The finite sum of error terms is bounded by
$4\,2^{-x^2/4}$.
\end{proof}

\begin{lemma}\label{lem:normalizer}
For $n\ge1$ and $0\le q\le n$,
\begin{equation}\label{eq:normalizer}
 B_n\ge2^{-32}B_{n-1},\qquad B_{n-q}/B_n\le2^{32q}.
\end{equation}
\end{lemma}
\begin{proof}
Take a DAG on $n-1$ labels, with $S$ sources. Choose one of the $n$ labels
for a new vertex, connect it to every old source, and choose arbitrary
outgoing edges to the other $n-1-S$ vertices. The new vertex is the unique
source. Deleting it recovers the input and its chosen edges, so this is an
injective count, after the remaining labels are identified in increasing
order with $[n-1]$. Hence
\[
 a_n\ge n2^{n-1}a_{n-1}\E[2^{-S}],\qquad
 B_n\ge B_{n-1}\E[2^{-S}].
\]
Lemma~\ref{lem:layer} gives $\Prob(S\le31)\ge1/2$ and thus
$\E[2^{-S}]\ge2^{-32}$. For $n=1$ the empty input has $S=0$.
Iteration proves the second inequality.
\end{proof}

Call $q$ a layer boundary if it is the sum of an initial sequence of layer
sizes; $n$ is always a boundary.
\begin{lemma}\label{lem:prefix}
For every exact forward graph $F$ on $[q]$, where $1\le q\le n$,
\begin{equation}\label{eq:boundary}
 \Prob(q\text{ is a boundary},F_q=F)\le2^{32q}\Q_q(F).
\end{equation}
Consequently, for every integer $L\ge32$, $\ell=\min(n,L)$, and every
event $A$ of forward graphs on $[\ell]$,
\begin{equation}\label{eq:prefix}
 \Prob(F_\ell\in A)
 \le2L\,2^{64L}\Q_\ell(A)+4L\,2^{-L^2/4}.
\end{equation}
\end{lemma}
\begin{proof}
An admissible complete rank graph $J$ has contiguous source-layer blocks of
sizes $h$. Each permutation from ranks to original labels gives a pair
consisting of a labeled DAG and its chosen intralayer orders. All $n!$ pairs
are distinct, including when the DAG has automorphisms, because the orders
are part of the pair. Each has probability $1/(a_n\prod_i h_i!)$. Therefore
\[
 \Prob(F_n=J)=\frac{2^{-\binom n2}}{B_n\prod_i h_i!}.
\]
If $q$ is a boundary, the source layers of $F$ are contiguous blocks;
their vector $g=(g_1,\ldots,g_r)$ is uniquely determined by $F$.
The source layers of the induced remainder are exactly the remaining
layers of the full graph. Cross edges are free except that every vertex of
the first remainder layer needs a parent in the last prefix layer. Summing
over remainder graphs and cross edges gives the exact identity
\begin{equation}\label{eq:exactprefix}
 \Prob(q\text{ boundary},F_q=F)
 =\frac{2^{-\binom q2}}{B_n\prod_{i=1}^r g_i!}
   \sum_{|h|=n-q}W(h)(1-2^{-g_r})^{h_1}.
\end{equation}
When $q=n$, the sum and its interface factor are both one. Otherwise the
sum is at most $B_{n-q}$. Lemma~\ref{lem:normalizer} proves
\eqref{eq:boundary}. If $F$ has noncontiguous source layers, the event is
empty; there is no additional sum over prefix vectors.

For \eqref{eq:prefix}, if $n<L$ use the terminal boundary $q=n$.
If $n\ge L$, take the first boundary $q\ge L$. When $q>2L$, the layer
containing rank $L$ has size at least $L$, and its index is at most $L$.
Lemma~\ref{lem:layer} and a union bound give probability at most
$4L\,2^{-L^2/4}$ for this exception. Otherwise sum
\eqref{eq:boundary} over $L\le q\le\min(2L,n)$, restricting each graph
to its first $L$ ranks. Under $\Q_q$ that restriction has law $\Q_L$.
There are at most $L+1\le2L$ boundaries to consider, and
$2^{32q}\le2^{64L}$.
\end{proof}

We call an edge non-collider if it belongs to no unshielded collider.
\begin{lemma}\label{lem:deep}
There are universal $c_1,C_1>0$ such that, for all integers $n,L\ge1$,
\begin{equation}\label{eq:deep}
 \Prob(\text{some non-collider edge has an endpoint of rank}>L)
 \le C_1e^{-c_1L}.
\end{equation}
\end{lemma}
\begin{proof}
For a possible larger endpoint $b\ge6400$, let $r=\lfloor b/100\rfloor$.
By Lemma~\ref{lem:layer}, the chance that one of the first $b$ layers has
size greater than $r$ is at most $4b\,2^{-r^2/4}$. Condition on ordered
layers satisfying this size cutoff. An edge $a\to b$, $a<b$, must join
distinct layers. Among the first $\lfloor b/2\rfloor$ ranks, discard those
in layers at distance at most one from either endpoint's layer.
At most $6r$ ranks are discarded, leaving at least $b/3$ witnesses.
Every layer meeting these earlier ranks has index at most $b$; a layer
after the layer of $b$ contains none of them.

For each remaining witness $u$, the bits $\{u,b\}$ and $\{u,a\}$ are
independent fair bits by Lemma~\ref{lem:enumeration}. These pairs are also
disjoint across witnesses. Since $u<b$, the first edge, if present,
points to $b$. With probability $1/4$ it is present and $\{u,a\}$ is
absent. If $a\to b$ is present, this gives $a\to b\leftarrow u$.
Thus the event that $a\to b$ exists and is non-collider has conditional
probability at most $(3/4)^{b/3}$. We have not conditioned on the existence
of $a\to b$, so adjacent-layer restrictions on that edge cause no problem.

A union bound over $a<b$ gives
\[
 b(3/4)^{b/3}+4b\,2^{-\lfloor b/100\rfloor^2/4}.
\]
Since $\lfloor b/100\rfloor\ge b/200$, this is at most $A e^{-cb}$
for universal positive constants. Sum over $b>L$, extending the sum to
infinity to remove dependence on $n$. Increasing $C_1$ covers $L<6400$.
\end{proof}

\section{Quadratic cost of the component conditions}
We first work under the independent-edge law $\Q$.
\begin{lemma}\label{lem:colliderfree}
There are universal $a,A>0$ such that, for every integer $r\ge2$,
\begin{equation}\label{eq:colliderfree}
 \Q_r(\text{no unshielded collider})\le e^{-ar^2+Ar}.
\end{equation}
\end{lemma}
\begin{proof}
For $r\ge16$ let $U$ be the first $\lfloor r/2\rfloor$ vertices, $V$ the
remaining vertices, and $k=\lfloor r/8\rfloor\ge r/16$. Reveal only the
edges within $U$. If the graph of missing skeleton edges on $U$ has a
matching of size $k$, choose one by a deterministic rule. For each $v\in V$,
absence of colliders requires that $v$ not receive both endpoints of any
matched pair. The parent bits are distinct and unrevealed. Hence the
conditional probability is at most
\[
 (3/4)^{k|V|}\le\exp[-\log(4/3)\,r^2/32].
\]
If no such matching exists, a maximum matching leaves at least
$\lfloor r/2\rfloor-2(k-1)\ge r/4-1$ unmatched vertices. No missing
edge joins two of them. They therefore form a clique of size at least $k$.
The probability of this alternative is at most
\[
 2^r2^{-\binom k2}
 \le\exp[-(\log2)r^2/512+(33/32)(\log2)r],
\]
where $\binom k2\ge r^2/512-r/32$ follows from $k\ge r/16$.
Adding the two alternatives proves \eqref{eq:colliderfree} with, for example,
\[
 a=\min\{\log(4/3)/32,\log2/512\},\qquad A=2\log2.
\]
For $2\le r<16$ the right side of \eqref{eq:colliderfree} exceeds one.
\end{proof}

Select $s$ vertices of a $\Q_\ell$ graph, partition them into blocks
$B_1,\ldots,B_k$ of sizes $r_i\ge2$, and fix a proposed topological order
of the blocks. Require each induced block to have no unshielded collider.
For $i<j$, require each vertex of $B_i$ to be a parent either of all
vertices of $B_j$ or of none, and forbid arrows from $B_j$ to $B_i$.
These are necessary conditions from Lemma~\ref{lem:components} for distinct
essential components, or for subsets of them.

\begin{lemma}\label{lem:partition}
For a fixed selected set, partition and block order, the probability of
these conditions is at most
\begin{equation}\label{eq:fixedpartition}
 \exp(-a_0s^2+As),\qquad a_0=\min\{a,\log2/4\}>0.
\end{equation}
Let $A_{m,\ell}$ be the event that some selected set of size
$m\le s\le2m$ admits such a partition and order. There is a universal
$d>0$ such that, for every sufficiently large integer $m$ and
$\ell\le L\le m^2$,
\begin{equation}\label{eq:partition}
 \Q_\ell(A_{m,\ell})\le e^{-dm^2}.
\end{equation}
\end{lemma}
\begin{proof}
For a fixed $u\in B_i$, $i<j$, at most two patterns on its $r_j$ pairs
with $B_j$ are allowed: all arrows $u\to B_j$, or all pairs absent.
If a rank in $B_j$ precedes $u$, the all-parent pattern is impossible,
which only decreases the count. The interblock probability is therefore
at most $2^{-r_i(r_j-1)}$. Since $r_j\ge2$,
\[
 \sum_{i<j}r_i(r_j-1)
 \ge\tfrac12\sum_{i<j}r_ir_j
 =\tfrac14\left(s^2-\sum_i r_i^2\right).
\]
Internal and interblock edge sets are disjoint. The induced law within
each selected block is $\Q_{r_i}$, regardless of gaps between its ranks.
Lemma~\ref{lem:colliderfree} thus bounds the full probability by
\[
 \exp\left[-a\sum_i r_i^2+As
       -\frac{\log2}{4}\left(s^2-\sum_i r_i^2\right)\right]
 \le e^{-a_0s^2+As}.
\]

There are at most $L^s$ choices of vertices, $s^s$ partitions and $s!$
block orders. For $m\le s\le2m$, $L\le m^2$, the logarithm of this
overcount is at most $s(4\log m+2\log2)$. Together with $As$, it is
at most $(a_0/2)s^2$ for all sufficiently large $m$. Sum the resulting
bound $e^{-a_0m^2/2}$ over at most $m+1$ values of $s$.
One may take $d=a_0/4$ after increasing the threshold. Values $s>\ell$
give empty events.
\end{proof}

\section{Proof of the tails and their scale}
\begin{proof}[Proof of Theorem~\ref{thm:vertices}]
Let $D=64\log2+1$, choose $\kappa=d/(4D)<1$, and put
$L=\lfloor\kappa m^2\rfloor$. For large enough $m$, we have
$32\le L\le m^2$ and $L\ge\kappa m^2/2$. If $n<m$ there is nothing
to prove. Otherwise, on $R\ge m$, choose $m$ vertices from one component
of size at least $m$, if such a component exists. If all components are
smaller, accumulate whole nontrivial components until the total first
reaches $m$. Their total size $s$ then lies in $[m,2m)$.
In either case the selected blocks have at least two vertices and inherit
the conditions of Lemma~\ref{lem:components}. They need not be essential
components of the induced prefix graph.

By Lemma~\ref{lem:components}, every vertex counted by $R$ is incident to
a non-collider edge. Outside the exception of Lemma~\ref{lem:deep}, all
such vertices have rank at most $L$. Thus the selection witnesses
$A_{m,\ell}$ in $F_\ell$, where $\ell=\min(n,L)$. Lemmas~\ref{lem:prefix}
and \ref{lem:partition} give
\begin{equation}\label{eq:combine}
 \Prob(R\ge m)
 \le C_1 e^{-c_1L}
       +2L\,2^{64L}e^{-dm^2}
       +4L\,2^{-L^2/4}.
\end{equation}
Since $\log(2L)\le L$, the middle term is at most
$\exp(DL-dm^2)\le e^{-3dm^2/4}$. The first term is at most
$C_1e^{-c_1\kappa m^2/2}$. The last is bounded by
$C'e^{-c'm^4}$, and hence by a Gaussian in $m$.
All constants are independent of $n$. Increasing $C$ covers the finitely
many smaller positive integers $m$.
\end{proof}

\begin{proof}[Proof of Theorem~\ref{thm:main} and Corollary~\ref{cor:target}]
Write $x=\log t$ and, for large enough $x$, put
\[
 m=\lfloor x/(2\log x)\rfloor.
\]
Then $m\ge x/(4\log x)$, while
\[
 \log(m!)\le m\log m\le m\log x\le x/2.
\]
In particular $m!\le t$. By \eqref{eq:factorial}, $|\M(G_n)|>t$
implies $R(G_n)>m$. Theorem~\ref{thm:vertices} now proves
\eqref{eq:main}, after changing its universal constant $c$.
For each fixed $\varepsilon\in(0,1)$,
$x^\varepsilon/(\log x)^2\to\infty$. Choose $X_\varepsilon$ above the
preceding thresholds so that $c x^\varepsilon/(\log x)^2\ge1$ for
all $x\ge X_\varepsilon$. With $t_\varepsilon=e^{X_\varepsilon}$ and
$C_\varepsilon=C$, \eqref{eq:main} gives Corollary~\ref{cor:target}.
\end{proof}

\begin{proposition}\label{prop:sharp}
The Gaussian scale in \eqref{eq:vertices} and the scale
$(\log t/\log\log t)^2$ in \eqref{eq:main} are sharp up to universal
constants in the exponent along complete-skeleton examples. In particular,
no constants $c,C>0$ give a uniform upper bound
$C\exp[-c(\log t)^2]$ for the class-size tail at all large $t$.
\end{proposition}
\begin{proof}
Every DAG has a topological order, so $a_n\le n!2^{\binom n2}$.
Exactly $n!$ DAGs have complete skeleton, one for each total order.
They belong to one class, of size $n!$, and have $R=n$ for $n\ge2$.
Consequently
\[
 \Prob(\text{complete skeleton})\ge2^{-\binom n2}.
\]
Also $n\log n-n+1\le\log(n!)\le n\log n$, by integrating the
increasing function $\log x$. For $t_n=n!-1$ we have
$\log t_n\sim n\log n$ and
$\log t_n/\log\log t_n\sim n$. The displayed lower bound thus has
the same exponential scales as \eqref{eq:vertices} and \eqref{eq:main}.
A proposed bound $C e^{-c(\log t)^2}$ at $t=t_n$ would imply
$c(\log t_n)^2\le\log C+(\log2)\binom n2$, which is impossible
after division by $n^2$.
\end{proof}

\paragraph{Finite check.}
A standalone exact-enumeration script accompanies the manuscript.
For four vertices it enumerates all 543 labeled DAGs and their 185 Markov
equivalence classes. It checks the source-layer weights, the exact prefix
identity \eqref{eq:exactprefix}, the common-parent condition and the
factorial bound using integer and rational arithmetic. These finite checks
are ancillary; no computation is used as a premise of the proofs above.

\begin{thebibliography}{9}
\bibitem{AMP}
S. A. Andersson, D. Madigan and M. D. Perlman,
\emph{A characterization of Markov equivalence classes for acyclic digraphs},
Annals of Statistics \textbf{25} (1997), no. 2, 505--541.
\href{https://doi.org/10.1214/aos/1031833662}{doi:10.1214/aos/1031833662}.
Primary-author version: University of Washington Technical Report 287 (1995),
\href{https://stat.uw.edu/sites/default/files/files/reports/1995/tr287.pdf}{full text}.
The locators used in this article refer to that report's printed pagination.

\bibitem{SS}
D. Schmid and A. Sly,
\emph{On the number and size of Markov equivalence classes of random
directed acyclic graphs},
\href{https://arxiv.org/abs/2209.04395v2}{arXiv:2209.04395v2} (2024),
revised September 28, 2024.
\end{thebibliography}

\paragraph{AI use.}
OpenAI Codex (GPT-6.1 Sol) was used to develop and check the argument,
verify source passages, implement the supplementary finite enumeration,
and prepare this manuscript. The manuscript has not been formally verified
or independently reviewed by a human. Independent mathematical review is
requested.
\end{document}
